\section{Detalle de todas las ideas exploradas}

A lo largo del proceso inicial de ideación, se generaron y documentaron diversas alternativas tecnológicas (13 ideas en total). Para estructurar este anexo de forma progresiva, las propuestas alternativas (excluyendo a \textit{SocioUnido}, que fue la idea finalmente seleccionada y detallada en el cuerpo principal del documento) se han clasificado en tres niveles según la profundidad del análisis realizado:

% ==========================================
% NIVEL 1
% ==========================================
\subsection{Nivel 1: Propuestas con planteo inicial}

Las siguientes siete propuestas fueron descartadas en la etapa temprana de validación.

\begin{itemize}
    \item \textbf{ScoreMaster (Promiedos 2.0)} 
    \newline \textit{Necesidad:} Las plataformas de seguimiento deportivo actuales presentan problemas críticos de rendimiento y una experiencia de usuario (UX) obsoleta. 
    \newline \textit{Solución:} Desarrollo de una plataforma web/móvil optimizada para datos en tiempo real mediante WebSockets, incorporando IA para generar de las jornadas deportivas.
    
    \item \textbf{AeroGuard}
    \newline \textit{Necesidad:} Riesgo latente de incursiones y colisiones en pistas de aeropuertos por errores de mando o vehículos no autorizados.
    \newline \textit{Solución:} Sistema de visión computacional y cámaras de alta definición que actúa como un \textbf{ojo digital} para auditar pistas y emitir alertas lumínicas autónomas ante alertas.
    
    \item \textbf{SkyDictate}
    \newline \textit{Necesidad:} Errores humanos de comunicación por radio entre la torre de control y los pilotos, derivando en quiebres de separación aérea.
    \newline \textit{Solución:} Software de auditoría en tiempo real que transcribe (conversión de voz a texto [Speech-to-Text]) y valida las instrucciones del controlador contra las normativas de altura, emitiendo alertas previas a la ejecución de maniobras peligrosas.

    \item \textbf{AutoCare}
    \newline \textit{Necesidad:} El mantenimiento automotriz suele ser puramente reactivo y desorganizado, reduciendo la vida útil del vehículo.
    \newline \textit{Solución:} Ecosistema preventivo que utiliza Inteligencia Artificial (reconocimiento óptico de caracteres [OCR]) para leer facturas de taller, digitalizar el historial del vehículo y cruzarlo con el manual del fabricante para predecir cuándo fallarán los repuestos.

    \item \textbf{NutriPlan}
    \newline \textit{Necesidad:} Frustración al planificar comidas que equilibren gustos, presupuesto, restricciones médicas y desgaste físico.
    \newline \textit{Solución:} Plataforma que cruza las preferencias del usuario, intolerancias y nivel de actividad, utilizando IA para generar un menú semanal hiper-personalizado y seguro.

    \item \textbf{ServiHogar}
    \newline \textit{Necesidad:} Dificultad histórica, especialmente en adultos mayores, para encontrar profesionales de oficios (plomeros, electricistas) confiables.
    \newline \textit{Solución:} Mercado digital (Marketplace) bidireccional con un enfoque de accesibilidad extrema y modelos de IA predictiva que estiman rangos de precios justos al describir el problema.

    \item \textbf{MarketScout}
    \newline \textit{Necesidad:} Tareas tediosas al comparar precios y promociones bancarias a través de múltiples supermercados online.
    \newline \textit{Solución:} Motor unificador de catálogos donde el usuario arma un único carrito y la IA (mediante \textit{extracción de datos [scraping]} y normalización) le indica en qué cadena resulta más barato efectuar la compra total o dividida.
\end{itemize}

\newpage

% ==========================================
% NIVEL 2
% ==========================================
\subsection{Nivel 2: Propuestas con análisis de contexto (5 C)}

\subsubsection{1. CryptoWealth: Gestión integral de activos digitales}

\begin{itemize}
    \item \textbf{Necesidad:} Fragmentación patrimonial al tener criptomonedas y acciones distribuidas en múltiples plataformas de intercambio (exchanges) y corredores de bolsa (brokers).
    \item \textbf{Solución:} Plataforma centralizada que consolida activos y utiliza PLN para hacer un \textit{análisis de sentimiento} de las noticias financieras, prediciendo movimientos del mercado.
\end{itemize}

% TABLAS 5C CRYPTOWEALTH
\begin{table}[h!]
\centering
\renewcommand{\arraystretch}{1.5}
\begin{tabularx}{\textwidth}{|X|}
\hline
\rowcolor{blue!20} \multicolumn{1}{|c|}{\textbf{COMPAÑÍA}} \\ \hline
\rowcolor{blue!10} \multicolumn{1}{|c|}{\textbf{Descripción}} \\ \hline
Empresa de tecnología financiera que ofrece una plataforma centralizada en la nube para gestión de activos. Incorpora IA basada en PLN para anticipar movimientos del mercado. \\ \hline
\rowcolor{blue!10} \multicolumn{1}{|c|}{\textbf{Visión}} \\ \hline
Convertirse en el estándar global de gestión patrimonial digital, eliminando la fragmentación entre cripto y finanzas tradicionales. \\ \hline
\rowcolor{blue!10} \multicolumn{1}{|c|}{\textbf{Misión}} \\ \hline
Brindando herramientas predictivas que permitan tomar decisiones informadas optimizando el rendimiento mediante consolidación unificada y segura. \\ \hline
\end{tabularx}
\end{table}

\begin{table}[h!]
\centering
\renewcommand{\arraystretch}{1.5}
\begin{tabularx}{\textwidth}{|X|}
\hline
\rowcolor{cyan!20} \multicolumn{1}{|c|}{\textbf{CLIENTES}} \\ \hline
\textbf{B2C (Empresa a Consumidor):} Inversores minoristas (retail), jóvenes nativos digitales e inversores avanzados con carteras diversificadas.\\
\textbf{E2E (Empresa a Empresa):} Asesores financieros, \textit{oficinas familiares (Family Offices)} y pequeños fondos de inversión que necesitan reportar múltiples clientes. \\ \hline
\end{tabularx}
\end{table}

\begin{table}[h!]
\centering
\renewcommand{\arraystretch}{1.5}
\begin{tabularx}{\textwidth}{|X|}
\hline
\rowcolor{green!20} \multicolumn{1}{|c|}{\textbf{COLABORADORES}} \\ \hline
Proveedores en la nube (\textit{Cloud}), interfaces de programación de aplicaciones (APIs) de cotización, agregadores de noticias (Bloomberg) y entidades regulatorias (CNV/UIF). \\ \hline
\end{tabularx}
\end{table}

\begin{table}[h!]
\centering
\renewcommand{\arraystretch}{1.5}
\begin{tabularx}{\textwidth}{|X|}
\hline
\rowcolor{orange!20} \multicolumn{1}{|c|}{\textbf{COMPETIDORES}} \\ \hline
Seguidores de portafolio (Portfolio trackers como CoinStats, Delta), plataformas propietarias de plataformas de intercambio, y el estado actual (\textit{status quo}): planillas de cálculo (Excel). \\ \hline
\end{tabularx}
\end{table}

\begin{table}[h!]
\centering
\renewcommand{\arraystretch}{1.5}
\begin{tabularx}{\textwidth}{|l|X|X|X|}
\hline
\rowcolor{red!20} \multicolumn{4}{|c|}{\textbf{CONTEXTO}} \\ \hline
\rowcolor{red!10} \textbf{Categoría} & \textbf{Corto plazo} & \textbf{Mediano plazo} & \textbf{Largo plazo} \\ \hline
\textbf{Económico} & \cellcolor{orange!50} & \cellcolor{green!50} & \cellcolor{green!50} \\ \hline
\textbf{Regulatorio} & \cellcolor{red!50} & \cellcolor{orange!50} & \cellcolor{green!50} \\ \hline
\textbf{Social / Cultural} & \cellcolor{green!50} & \cellcolor{green!50} & \cellcolor{green!50} \\ \hline
\textbf{Tecnológico} & \cellcolor{green!50} & \cellcolor{green!50} & \cellcolor{green!50} \\ \hline
\textbf{Político / Institucional} & \cellcolor{orange!50} & \cellcolor{orange!50} & \cellcolor{green!50} \\ \hline
\end{tabularx}
\end{table}

\newpage

\subsubsection{2. MedTrack: Digitalización y seguimiento médico}
\begin{itemize}
    \item \textbf{Necesidad:} Fragmentación en clínicas privadas, pérdida de historiales e ineficiencia en agendas.
    \item \textbf{Solución:} Sistema unificado de historias clínicas en la nube con algoritmos de optimización para reducir el ausentismo de pacientes.
\end{itemize}

% TABLAS 5C MEDTRACK
\begin{table}[h!]
\centering
\renewcommand{\arraystretch}{1.5}
\begin{tabularx}{\textwidth}{|X|}
\hline
\rowcolor{blue!20} \multicolumn{1}{|c|}{\textbf{COMPAÑÍA}} \\ \hline
\rowcolor{blue!10} \multicolumn{1}{|c|}{\textbf{Descripción}} \\ \hline
Empresa de tecnología de la salud (\textit{HealthTech}) enfocada en resolver la fragmentación clínica digitalizando historias médicas y optimizando agendas mediante IA. \\ \hline
\rowcolor{blue!10} \multicolumn{1}{|c|}{\textbf{Visión}} \\ \hline
Ser el estándar de interoperabilidad clínica garantizando una atención fluida y sin burocracia. \\ \hline
\rowcolor{blue!10} \multicolumn{1}{|c|}{\textbf{Misión}} \\ \hline
Brindando herramientas tecnológicas que eliminen la pérdida de datos y el ausentismo, mejorando drásticamente la calidad de atención \\ \hline
\end{tabularx}
\end{table}

\begin{table}[h!]
\centering
\renewcommand{\arraystretch}{1.5}
\begin{tabularx}{\textwidth}{|X|}
\hline
\rowcolor{cyan!20} \multicolumn{1}{|c|}{\textbf{CLIENTES}} \\ \hline
\textbf{E2E:} Clínicas, sanatorios, centros de diagnóstico y obras sociales.\\
\textbf{B2C:} Médicos especialistas independientes y pequeños consultorios. \\ \hline
\end{tabularx}
\end{table}

\begin{table}[h!]
\centering
\renewcommand{\arraystretch}{1.5}
\begin{tabularx}{\textwidth}{|X|}
\hline
\rowcolor{green!20} \multicolumn{1}{|c|}{\textbf{COLABORADORES}} \\ \hline
Proveedores en la nube (\textit{Cloud} de alta seguridad), integradores de pago, Obras Sociales, y Ministerios de Salud para firmas electrónicas. \\ \hline
\end{tabularx}
\end{table}

\begin{table}[h!]
\centering
\renewcommand{\arraystretch}{1.5}
\begin{tabularx}{\textwidth}{|X|}
\hline
\rowcolor{orange!20} \multicolumn{1}{|c|}{\textbf{COMPETIDORES}} \\ \hline
Software médico directo (Doctoralia, Omnia), sistemas cerrados de grandes hospitales, y herramientas no especializadas (Agendas papel, WhatsApp). \\ \hline
\end{tabularx}
\end{table}

\begin{table}[h!]
\centering
\renewcommand{\arraystretch}{1.5}
\begin{tabularx}{\textwidth}{|l|X|X|X|}
\hline
\rowcolor{red!20} \multicolumn{4}{|c|}{\textbf{CONTEXTO}} \\ \hline
\rowcolor{red!10} \textbf{Categoría} & \textbf{Corto plazo} & \textbf{Mediano plazo} & \textbf{Largo plazo} \\ \hline
\textbf{Económico} & \cellcolor{red!50} & \cellcolor{orange!50} & \cellcolor{green!50} \\ \hline
\textbf{Regulatorio} & \cellcolor{orange!50} & \cellcolor{green!50} & \cellcolor{green!50} \\ \hline
\textbf{Social / Cultural} & \cellcolor{orange!50} & \cellcolor{green!50} & \cellcolor{green!50} \\ \hline
\textbf{Tecnológico} & \cellcolor{green!50} & \cellcolor{green!50} & \cellcolor{green!50} \\ \hline
\textbf{Político / Institucional} & \cellcolor{green!50} & \cellcolor{green!50} & \cellcolor{green!50} \\ \hline
\end{tabularx}
\end{table}

\newpage

\subsubsection{3. SmartPantry: Inventario gastronómico inteligente}
\begin{itemize}
    \item \textbf{Necesidad:} Desperdicio constante de insumos perecederos en restaurantes y hogares debido al mal manejo del stock y caducidad.
    \item \textbf{Solución:} Motor de inventario dinámico que cruza fechas de vencimiento con modelos de lenguaje de gran tamaño (LLMs) para sugerir recetas y evitar mermas financieras.
\end{itemize}

% TABLAS 5C SMARTPANTRY
\begin{table}[h!]
\centering
\renewcommand{\arraystretch}{1.5}
\begin{tabularx}{\textwidth}{|X|}
\hline
\rowcolor{blue!20} \multicolumn{1}{|c|}{\textbf{COMPAÑÍA}} \\ \hline
\rowcolor{blue!10} \multicolumn{1}{|c|}{\textbf{Descripción}} \\ \hline
Empresa emergente (\textit{Startup}) de tecnología alimentaria (\textit{FoodTech}) que utiliza IA para priorizar el consumo de alimentos en base a fechas de caducidad, minimizando el desperdicio. \\ \hline
\rowcolor{blue!10} \multicolumn{1}{|c|}{\textbf{Visión}} \\ \hline
Liderar la gestión inteligente de inventarios gastronómicos en LATAM para hacer la industria más sostenible. \\ \hline
\rowcolor{blue!10} \multicolumn{1}{|c|}{\textbf{Misión}} \\ \hline
Automatizando el stock y sugiriendo recetas para optimizar compras y rentabilidad. \\ \hline
\end{tabularx}
\end{table}

\begin{table}[h!]
\centering
\renewcommand{\arraystretch}{1.5}
\begin{tabularx}{\textwidth}{|X|}
\hline
\rowcolor{cyan!20} \multicolumn{1}{|c|}{\textbf{CLIENTES}} \\ \hline
\textbf{E2E:} Restaurantes, bares y servicios de viandas. \\
\textbf{B2C:} Familias e individuos orientados a reducir la huella de desperdicio y optimizar su alacena. \\ \hline
\end{tabularx}
\end{table}

\begin{table}[h!]
\centering
\renewcommand{\arraystretch}{1.5}
\begin{tabularx}{\textwidth}{|X|}
\hline
\rowcolor{green!20} \multicolumn{1}{|c|}{\textbf{COLABORADORES}} \\ \hline
Proveedores de alimentos, escuelas de gastronomía (para validar recetas de la IA) y ONGs medioambientales. \\ \hline
\end{tabularx}
\end{table}

\begin{table}[h!]
\centering
\renewcommand{\arraystretch}{1.5}
\begin{tabularx}{\textwidth}{|X|}
\hline
\rowcolor{orange!20} \multicolumn{1}{|c|}{\textbf{COMPETIDORES}} \\ \hline
Sistemas de planificación de recursos empresariales (ERP) genéricos, aplicaciones (apps) de recetas estándar, e inventarios en Excel. \\ \hline
\end{tabularx}
\end{table}

\begin{table}[h!]
\centering
\renewcommand{\arraystretch}{1.5}
\begin{tabularx}{\textwidth}{|l|X|X|X|}
\hline
\rowcolor{red!20} \multicolumn{4}{|c|}{\textbf{CONTEXTO}} \\ \hline
\rowcolor{red!10} \textbf{Categoría} & \textbf{Corto plazo} & \textbf{Mediano plazo} & \textbf{Largo plazo} \\ \hline
\textbf{Económico} & \cellcolor{green!50} & \cellcolor{green!50} & \cellcolor{green!50} \\ \hline
\textbf{Regulatorio} & \cellcolor{green!50} & \cellcolor{green!50} & \cellcolor{green!50} \\ \hline
\textbf{Social / Cultural} & \cellcolor{orange!50} & \cellcolor{orange!50} & \cellcolor{green!50} \\ \hline
\textbf{Tecnológico} & \cellcolor{green!50} & \cellcolor{green!50} & \cellcolor{green!50} \\ \hline
\textbf{Político / Institucional} & \cellcolor{green!50} & \cellcolor{green!50} & \cellcolor{green!50} \\ \hline
\end{tabularx}
\end{table}

\newpage

% ==========================================
% NIVEL 3
% ==========================================
\subsection{Nivel 3: Propuestas con análisis de negocio y marketing (5 C y 4 P)}

\subsubsection{1. DevQuest: Gamificación de la inducción en TI}

\begin{itemize}
    \item \textbf{Necesidad:} Cuellos de botella formativos al ingresar un nuevo desarrollador a una empresa, ralentizando la productividad de los desarrolladores experimentados (Seniors).
    \item \textbf{Solución:} Plataforma de software como servicio (SaaS) que transforma la inducción en un \textbf{árbol de habilidades} gamificado, revisando código estático con IA y resolviendo dudas del código legado mediante un bot de generación aumentada por recuperación (RAG).
\end{itemize}

% TABLAS 5C DEVQUEST
\begin{table}[h!]
\centering
\renewcommand{\arraystretch}{1.5}
\begin{tabularx}{\textwidth}{|X|}
\hline
\rowcolor{blue!20} \multicolumn{1}{|c|}{\textbf{COMPAÑÍA}} \\ \hline
\rowcolor{blue!10} \multicolumn{1}{|c|}{\textbf{Descripción}} \\ \hline
Empresa emergente de tecnología aplicada a recursos humanos (HR-Tech) que reduce la carga sobre los líderes técnicos acelerando la inducción mediante gamificación y generación aumentada por recuperación (RAG). \\ \hline
\rowcolor{blue!10} \multicolumn{1}{|c|}{\textbf{Visión}} \\ \hline
Erradicar la fricción en la integración de talento permitiendo productividad autónoma desde el primer día en cualquier conjunto de tecnologías (stack). \\ \hline
\rowcolor{blue!10} \multicolumn{1}{|c|}{\textbf{Misión}} \\ \hline
Digitalizando la inducción (\textit{onboarding}), liberando tiempo estratégico de las gerencias. \\ \hline
\end{tabularx}
\end{table}

\begin{table}[h!]
\centering
\renewcommand{\arraystretch}{1.5}
\begin{tabularx}{\textwidth}{|X|}
\hline
\rowcolor{cyan!20} \multicolumn{1}{|c|}{\textbf{CLIENTES}} \\ \hline
Empresas de software (fábricas de software [\textit{Software Factories}], empresas emergentes [Startups]), directores de tecnología (CTOs) y líderes técnicos (gerentes de ingeniería). \\ \hline
\end{tabularx}
\end{table}

\begin{table}[h!]
\centering
\renewcommand{\arraystretch}{1.5}
\begin{tabularx}{\textwidth}{|X|}
\hline
\rowcolor{green!20} \multicolumn{1}{|c|}{\textbf{COLABORADORES}} \\ \hline
Comunidades de TI, agencias de reclutamiento (recruiting) y la urgencia del mercado por adquirir talento especializado en nuevas arquitecturas. \\ \hline
\end{tabularx}
\end{table}

\begin{table}[h!]
\centering
\renewcommand{\arraystretch}{1.5}
\begin{tabularx}{\textwidth}{|X|}
\hline
\rowcolor{orange!20} \multicolumn{1}{|c|}{\textbf{COMPETIDORES}} \\ \hline
Burocracia interna (acompañamiento [\textit{Shadowing}] manual en videollamadas), herramientas de recursos humanos genéricas (BambooHR), y la penetración de modelos de IA generativa. \\ \hline
\end{tabularx}
\end{table}

\begin{table}[h!]
\centering
\renewcommand{\arraystretch}{1.5}
\begin{tabularx}{\textwidth}{|l|X|X|X|}
\hline
\rowcolor{red!20} \multicolumn{4}{|c|}{\textbf{CONTEXTO}} \\ \hline
\rowcolor{red!10} \textbf{Categoría} & \textbf{Corto plazo} & \textbf{Mediano plazo} & \textbf{Largo plazo} \\ \hline
\textbf{Económico} & \cellcolor{green!50} & \cellcolor{green!50} & \cellcolor{green!50} \\ \hline
\textbf{Regulatorio} & \cellcolor{green!50} & \cellcolor{green!50} & \cellcolor{green!50} \\ \hline
\textbf{Social / Cultural} & \cellcolor{green!50} & \cellcolor{green!50} & \cellcolor{green!50} \\ \hline
\textbf{Tecnológico} & \cellcolor{green!50} & \cellcolor{green!50} & \cellcolor{green!50} \\ \hline
\textbf{Político / Ambiental} & \cellcolor{green!50} & \cellcolor{green!50} & \cellcolor{green!50} \\ \hline
\end{tabularx}
\end{table}

\newpage

\textbf{Análisis de las 4 P (DevQuest)}
\begin{itemize}
    \item \textbf{Producto:} Sistema de \textit{software como servicio (SaaS) estándar} (con árboles pre-creados para lenguajes populares) y \textit{proyectos empresariales (Enterprise)} a medida. Como ecosistema se planificó DevQuest Analítica (Analytics) y alianzas de certificación.
    \item \textbf{Plaza:} Distribución 100\% de empresa a empresa (E2E) enfocada inicialmente en fábricas de software (Software Factories) locales (50 a 200 empleados) por su alta rotación operativa.
    \item \textbf{Precio:} \textit{Fijación de precios (Pricing)} por valor (ahorro de horas de desarrolladores experimentados [seniors]). Suscripciones recurrentes (inicial/crecimiento [Starter/Growth]) y un cobro de capital por desarrollo en integraciones complejas (\textit{personalizado para empresas}).
    \item \textbf{Promoción:} Tracción temprana mediante el boca a boca (\textit{Word of Mouth}) en comunidades técnicas, seguido de mercadotecnia pagada de empresa a empresa (marketing pago E2E mediante anuncios de LinkedIn [LinkedIn Ads]) apuntando a directores de tecnología (CTOs).
\end{itemize}

\subsubsection{2. CuentasClaras: Gestión integral y predictiva del hogar}
\begin{itemize}
    \item \textbf{Necesidad:} Conflictos por división de dinero entre compañeros de cuarto (\textit{roommates}) y desorganización personal ante deudas, vencimientos e inflación.
    \item \textbf{Solución:} Billetera colaborativa que automatiza los pagos fijos, alerta sobre vencimientos y usa modelos predictivos para proyectar los presupuestos del mes ajustados por inflación.
\end{itemize}

% TABLAS 5C CUENTAS CLARAS
\begin{table}[h!]
\centering
\renewcommand{\arraystretch}{1.5}
\begin{tabularx}{\textwidth}{|X|}
\hline
\rowcolor{blue!20} \multicolumn{1}{|c|}{\textbf{COMPAÑÍA}} \\ \hline
\rowcolor{blue!10} \multicolumn{1}{|c|}{\textbf{Descripción}} \\ \hline
Empresa emergente de tecnología financiera (Startup Fintech) enfocada en automatizar el pago de servicios, la división de gastos y la adaptación presupuestaria. \\ \hline
\rowcolor{blue!10} \multicolumn{1}{|c|}{\textbf{Visión}} \\ \hline
Transformar la economía doméstica en una experiencia invisible y completamente libre de conflictos sociales. \\ \hline
\rowcolor{blue!10} \multicolumn{1}{|c|}{\textbf{Misión}} \\ \hline
Centralizando finanzas y eliminando la carga mental de la organización económica. \\ \hline
\end{tabularx}
\end{table}

\begin{table}[h!]
\centering
\renewcommand{\arraystretch}{1.5}
\begin{tabularx}{\textwidth}{|X|}
\hline
\rowcolor{cyan!20} \multicolumn{1}{|c|}{\textbf{CLIENTES}} \\ \hline
Familias consolidadas, parejas convivientes, estudiantes universitarios (compañeros de cuarto o \textit{roommates}) y gestiones individuales. \\ \hline
\end{tabularx}
\end{table}

\begin{table}[h!]
\centering
\renewcommand{\arraystretch}{1.5}
\begin{tabularx}{\textwidth}{|X|}
\hline
\rowcolor{green!20} \multicolumn{1}{|c|}{\textbf{COLABORADORES}} \\ \hline
Pasarelas de pagos, entidades gubernamentales (datos INDEC/inflación), y redes de publicidad digital (Google AdMob). \\ \hline
\end{tabularx}
\end{table}

\begin{table}[h!]
\centering
\renewcommand{\arraystretch}{1.5}
\begin{tabularx}{\textwidth}{|X|}
\hline
\rowcolor{orange!20} \multicolumn{1}{|c|}{\textbf{COMPETIDORES}} \\ \hline
Splitwise, planillas de Excel, economía informal (manejo en efectivo puro) y módulos limitados de aplicaciones (apps) bancarias tradicionales. \\ \hline
\end{tabularx}
\end{table}

\newpage

\begin{table}[h!]
\centering
\renewcommand{\arraystretch}{1.5}
\begin{tabularx}{\textwidth}{|l|X|X|X|}
\hline
\rowcolor{red!20} \multicolumn{4}{|c|}{\textbf{CONTEXTO}} \\ \hline
\rowcolor{red!10} \textbf{Categoría} & \textbf{Corto plazo} & \textbf{Mediano plazo} & \textbf{Largo plazo} \\ \hline
\textbf{Económico} & \cellcolor{green!50} & \cellcolor{orange!50} & \cellcolor{green!50} \\ \hline
\textbf{Regulatorio} & \cellcolor{green!50} & \cellcolor{green!50} & \cellcolor{green!50} \\ \hline
\textbf{Social / Cultural} & \cellcolor{orange!50} & \cellcolor{green!50} & \cellcolor{green!50} \\ \hline
\textbf{Tecnológico} & \cellcolor{green!50} & \cellcolor{green!50} & \cellcolor{green!50} \\ \hline
\textbf{Político / Ambiental} & \cellcolor{green!50} & \cellcolor{orange!50} & \cellcolor{green!50} \\ \hline
\end{tabularx}
\end{table}

\textbf{Análisis de las 4 P (CuentasClaras)}

\begin{itemize}
    \item \textbf{Producto:} Plataforma dividida en un \textit{nivel gratuito con opciones de pago (Tier Freemium)} con anuncios y una variante \textit{de pago (Premium)} (individual y hogar) con integraciones completas, predicción inflacionaria por IA y reconocimiento óptico de caracteres (OCR) para despiece de recibos (tickets).
    \item \textbf{Plaza:} Tiendas de aplicaciones móviles (App Store/Google Play). Para el dimensionamiento, se estimó capturar un mercado disponible servible (SAM) de compradores inteligentes (\textit{smart-shoppers}) con una meta inicial de 47.000 usuarios activos.
    \item \textbf{Precio:} Modelo de micropagos de alta penetración (\$1 a \$3 dólares [USD]), sustentado paralelamente por monetización publicitaria (ingreso promedio por usuario gratuito [ARPU Free]) y futuras comisiones por colocación de empresa a empresa (E2E) de seguros.
    \item \textbf{Promoción:} Estrategia basada en el \textbf{efecto red} orgánico (los usuarios deben invitar a sus convivientes), apoyada por optimización para tiendas de aplicaciones (ASO) y *mercadotecnia de influenciadores (marketing influencers)* de educación financiera.
\end{itemize}

\newpage

\subsection{Proceso de decisión de proyecto final en base a las 5C y las 4P}

Para determinar la viabilidad y seleccionar la idea de negocio definitiva, se aplicó un filtro analítico estructurado en dos etapas. En la primera instancia, se evaluaron las seis propuestas preseleccionadas utilizando la matriz de las \textbf{5 C} para identificar barreras macroeconómicas, regulatorias y competitivas. En la segunda instancia, los tres proyectos finalistas que superaron el primer filtro fueron sometidos a un análisis de penetración y rentabilidad mediante la mezcla de mercadotecnia (\textbf{4 P}).

Para facilitar la lectura de este proceso de descarte, se emplea una matriz de calor (semáforo) donde el color \textbf{\colorbox{green!50}{Verde}} indica un entorno favorable o de alto valor, el \textbf{\colorbox{yellow!50}{Amarillo}} representa un riesgo intermedio o moderado, y el \textbf{\colorbox{red!50}{Rojo}} denota un escenario crítico o desfavorable.

\subsubsection{Primera etapa de evaluación: Matriz de las 5 C}

\begin{table}[h!]
\centering
\renewcommand{\arraystretch}{1.5}
\begin{tabularx}{\textwidth}{|l|X|X|X|X|X|}
\hline
\rowcolor{gray!20} \textbf{Proyecto} & \textbf{Compañía} & \textbf{Clientes} & \textbf{Colabor.} & \textbf{Competi.} & \textbf{Contexto} \\ \hline
\textbf{CryptoWealth} & \cellcolor{yellow!50}Intermedio & \cellcolor{yellow!50}Intermedio & \cellcolor{green!50}Favorable & \cellcolor{red!50}Crítico & \cellcolor{red!50}Crítico \\ \hline
\textbf{MedTrack} & \cellcolor{green!50}Favorable & \cellcolor{yellow!50}Intermedio & \cellcolor{yellow!50}Intermedio & \cellcolor{red!50}Crítico & \cellcolor{red!50}Crítico \\ \hline
\textbf{SmartPantry} & \cellcolor{yellow!50}Intermedio & \cellcolor{yellow!50}Intermedio & \cellcolor{yellow!50}Intermedio & \cellcolor{yellow!50}Intermedio & \cellcolor{yellow!50}Intermedio \\ \hline
\textbf{DevQuest} & \cellcolor{green!50}Favorable & \cellcolor{green!50}Favorable & \cellcolor{green!50}Favorable & \cellcolor{yellow!50}Intermedio & \cellcolor{green!50}Favorable \\ \hline
\textbf{CuentasClaras} & \cellcolor{green!50}Favorable & \cellcolor{green!50}Favorable & \cellcolor{yellow!50}Intermedio & \cellcolor{yellow!50}Intermedio & \cellcolor{green!50}Favorable \\ \hline
\textbf{SocioUnido} & \cellcolor{green!50}Favorable & \cellcolor{green!50}Favorable & \cellcolor{green!50}Favorable & \cellcolor{green!50}Favorable & \cellcolor{green!50}Favorable \\ \hline
\end{tabularx}
\caption{Matriz de calor de las 5 C para la preselección de proyectos.}
\label{tab:semaforo_5c}
\end{table}

\textbf{Justificación del primer filtrado:}
Como se observa en el Cuadro \ref{tab:semaforo_5c}, las primeras tres propuestas fueron descartadas por presentar riesgos inasumibles en el corto plazo. \textit{CryptoWealth} y \textit{MedTrack} evidenciaron un entorno altamente crítico (\colorbox{red!50}{Rojo}) en las variables de \textbf{Contexto} y \textbf{Competidores}; la primera por la alta incertidumbre regulatoria financiera, y la segunda por la severa burocracia institucional de las obras sociales y la competencia de sistemas cerrados hospitalarios. Por su parte, \textit{SmartPantry} no presentó bloqueos críticos, pero su evaluación resultó íntegramente intermedia (\colorbox{yellow!50}{Amarillo}), careciendo de un diferencial sólido para traccionar a usuarios habituados a la gestión gastronómica analógica.

En contraposición, \textbf{DevQuest}, \textbf{CuentasClaras} y \textbf{SocioUnido} lograron un perfil marcadamente verde, demostrando tracción, viabilidad técnica y mercados accesibles, lo que justificó su avance hacia la evaluación comercial detallada.

\newpage

\subsubsection{Segunda etapa de evaluación: Matriz de las 4 P (Proyectos finalistas)}

\begin{table}[h!]
\centering
\renewcommand{\arraystretch}{1.5}
\begin{tabularx}{\textwidth}{|l|X|X|X|X|}
\hline
\rowcolor{gray!20} \textbf{Proyecto finalista} & \textbf{Producto} & \textbf{Plaza} & \textbf{Precio} & \textbf{Promoción} \\ \hline
\textbf{DevQuest} & \cellcolor{green!50}Favorable & \cellcolor{yellow!50}Intermedio & \cellcolor{green!50}Favorable & \cellcolor{yellow!50}Intermedio \\ \hline
\textbf{CuentasClaras} & \cellcolor{green!50}Favorable & \cellcolor{green!50}Favorable & \cellcolor{yellow!50}Intermedio & \cellcolor{yellow!50}Intermedio \\ \hline
\textbf{SocioUnido} & \cellcolor{green!50}Favorable & \cellcolor{green!50}Favorable & \cellcolor{green!50}Favorable & \cellcolor{green!50}Favorable \\ \hline
\end{tabularx}
\caption{Matriz de calor de las 4 P para la selección del proyecto definitivo.}
\label{tab:semaforo_4p}
\end{table}

\textbf{Justificación de la decisión final:}
El análisis de las 4 P (Cuadro \ref{tab:semaforo_4p}) permitió exponer las vulnerabilidades comerciales de los dos primeros finalistas, decantando la elección unánime hacia SocioUnido:

\begin{itemize}
    \item \textbf{DevQuest:} Si bien posee un producto innovador y altamente valorado, su \textbf{Plaza} y \textbf{Promoción} sufren una penalización (\colorbox{yellow!50}{Amarillo}). El mercado E2E de recursos humanos en empresas de TI se encuentra saturado y los ciclos de venta (apuntando a CTOs) son extremadamente largos y complejos para una empresa emergente.
    \item \textbf{CuentasClaras:} Presenta una excelente Plaza (distribución digital masiva), pero su debilidad radica en su modelo de \textbf{Precio} y \textbf{Promoción} (\colorbox{yellow!50}{Amarillo}). Depender de micropagos B2C y publicidad dentro de la aplicación genera márgenes de ganancia mínimos, requiriendo un volumen de usuarios extraordinariamente alto para alcanzar el punto de equilibrio, lo cual es riesgoso en un mercado altamente competitivo.
    \item \textbf{SocioUnido (Selección Definitiva):} Logró la validación máxima (\colorbox{green!50}{Verde}) en todos sus ejes. Presenta un \textbf{Producto} que resuelve dolores financieros críticos; una \textbf{Plaza} virgen e identificada (categorías de ascenso); un modelo de \textbf{Precio} E2E mixto que garantiza rentabilidad sustentable (cobro fijo mensual más variable por transacción); y una estrategia de \textbf{Promoción} altamente efectiva basada en el efecto contagio y auditorías dirigenciales directas.
\end{itemize}

\textbf{Conclusión:} En base a la solidez ininterrumpida demostrada tanto en el análisis macro-contextual (5 C) como en su viabilidad de penetración y rentabilidad comercial (4 P), el equipo determinó que \textbf{SocioUnido} representa la propuesta más robusta, escalable y acorde a las exigencias integrales requeridas para el Trabajo Profesional de Ingeniería en Informática.
